% 第3章 3.2.1：l^2 上 A(x_n)=(x_n/n) 的谱与预解算子范数。
% (a) 特征值 1/n（n=1..40 画实心点）聚集到 0；0 属于谱但不是特征值（空心点）。
% (b) 实参数 lambda 处 ||(A-lambda I)^{-1}||=1/dist(lambda,sigma(A))，对数纵轴，
%     数据由同目录 generate.py 生成（resolvent.dat，超过 1000 的值截为 1000，纵轴上限 120）。
\begin{tikzpicture}[font=\figurefont\footnotesize]
\begin{axis}[name=A,width=96mm,height=24mm,
  xmin=-0.3,xmax=1.3,ymin=-1,ymax=1,
  axis x line=middle,axis y line=none,clip=false,
  xtick={0,0.25,0.5,1},xticklabels={$0$,$1/4$,$1/2$,$1$},
  x axis line style={draw=FigureStroke,line width=.6pt,-{Stealth[length=1.6mm]}},
  tick label style={font=\figurefont\scriptsize,text=FigureInk},
  xlabel={$\lambda$},xlabel style={at={(axis description cs:1,0.5)},anchor=west},
  title={(a) $\sigma(A)=\{1/n:n\geq1\}\cup\{0\}$},title style={font=\figurefont\footnotesize}]
\addplot[only marks,mark=*,mark size=1.3pt,FigureBlue,samples at={1,...,40}] ({1/x},{0});
\addplot[only marks,mark=o,mark size=2pt,FigureRed,thick] coordinates {(0,0)};
\node[font=\FigureNoteFont\scriptsize,text=FigureRed,anchor=south] at (axis cs:-0.05,0.25) {非特征值};
\node[font=\FigureNoteFont\scriptsize,text=FigureBlue,anchor=south] at (axis cs:0.75,0.25) {特征值$1/n$};
\end{axis}
\begin{axis}[at={(A.south)},anchor=north,yshift=-12mm,width=96mm,height=42mm,
  xmin=-0.3,xmax=1.3,ymode=log,ymin=0.8,ymax=120,
  axis x line=bottom,axis y line=left,
  xtick={-0.25,0,0.25,0.5,0.75,1,1.25},
  xticklabel style={/pgf/number format/fixed},
  ytick={1,10,100},
  xlabel={$\lambda$},ylabel={$\|(A-\lambda I)^{-1}\|$},
  xlabel style={font=\figurefont\scriptsize},ylabel style={font=\figurefont\scriptsize},
  tick label style={font=\figurefont\scriptsize,text=FigureInk},
  ymajorgrids,grid style={draw=FigureGrid,line width=.4pt},
  title={(b) 预解算子范数在谱附近无界},title style={font=\figurefont\footnotesize}]
\addplot[FigureBlue,line width=.9pt] table {figures/02-abstract-mathematics/03-function-spaces/tikz/c3-spectrum-resolvent/resolvent.dat};
\end{axis}
\end{tikzpicture}
